%% kapa: exam=1 round=32 year=2021 period=1 form=A subject=economics subject_ko=경제학원론 seq=9 number=49
%% source: Q-Net 감정평가사 기출문제 (cst00309) artlSeq=5208397 file=extracted/1008/5208397/A형/2021년도 제32회 감정평가사 1차 1교시 A형.pdf
%% conversion: kordoc 4.24.51 -> tools/gen1.py (auto); review=proofread-sonnet; flags=none
\begin{question}{49}
후생경제이론에 관한 설명으로 옳은 것은?
\begin{choices}
  \item 파레토(Pareto) 효율적인 상태는 파레토 개선이 가능한 상태를 뜻한다.
  \item 제2정리는 모든 사람의 선호가 오목성을 가지면 파레토 효율적인 배분은 일반경쟁 균형이 된다는 것이다.
  \item 제1정리는 모든 소비자의 선호체계가 약 단조성을 갖고 외부성이 존재하면 일반경쟁 균형의 배분은 파레토 효율적이라는 것이다.
  \item 제1정리는 완전경쟁시장 하에서 사익과 공익은 서로 상충된다는 것이다.
  \item 제1정리는 아담 스미스(A. Smith)의 ‘보이지 않는 손’의 역할을 이론적으로 뒷받침 해주는 것이다.
\end{choices}
\end{question}
\answer{5}
